\chapter{Discussion}
\label{ch:discussion}

The real-data undersampling approach and the comparison of three importance
methods across three models are the central contributions of this thesis. The
main result of that comparison is that the methods diverge, but they seem to do
so in an anticipated pattern, each showing the strengths and weaknesses that
were laid out throughout this thesis. The logistic regression is the simplest case. 
Because the $\ell_1$ penalty
already thinned most of the redundancy during training, all three methods agree
on the Debt ratio. This clean picture, however, comes from the weakest of the
three models. Whether a linear model can represent the full complexity of the
data remains open. The two ensembles predict better, but they receive the 
redundancy intact, since all $95$ features were deliberately kept for the integrity of the data and for
comparability with \textcite{wangliu2021}. Here the methods drift apart in
recognizable ways. The Gini ranking of the random forest carries three closely
related earnings measures on places one, three and four, which fits its described
tendency to split credit between near-identical features. Permutation drops
Persistent EPS entirely in both ensembles although the other methods rank it
first or second, which fits its described tendency to hide features with close
substitutes. SHAP, in turn, keeps the redundant families visible in every model,
which fits its described tendency to share credit between them. The redundancy
raised at the beginning therefore did not harm prediction, but it shaped
interpretation, as the disagreements between the methods concentrate largely
where features have near-identical partners. This consistency between what the
methods are known to do and what they visibly did here is the main insight of
the comparison. On its own each of these lists looks like the answer. 
Only the comparison shows that it is one of several, shaped by the model 
and the method behind it.

Read economically, and strictly as association rather than cause, the results
still tell one consistent story. Every model places a leverage measure at or
near the top, the Debt ratio for the logistic regression and Borrowing
dependency for the ensembles, so firms financed largely through borrowed capital
are the ones the models flag, which is also what one would expect economically.
At the same time the single most important feature is model-dependent, as the
random forest, nearly as strong as gradient boosting, ranks Persistent EPS first
under SHAP. Among the studies reviewed here, \textcite{sinap2025} is the only
other one applying SHAP to this dataset and reports Net Income to Total Assets,
in substance a return-on-assets measure, as the top feature of a stacking
classifier trained on SMOTE-balanced data and is followed by the Debt ratio. It is
notable that this list resembles the ranking of the linear model more than
those of the better-performing ensembles, where the Debt ratio sits far from
the top. Comparing these answers does not make one of them wrong. It shows that
a ranking depends on the whole pipeline that produces it, which supports
reading such lists with care.

The undersampling approach deserves the same two-sided reading. On the positive
side it achieves what it was chosen for. Every firm the models learned from
actually existed and the reported scores are measured on untouched test folds
with the natural $3.2\%$ bankruptcy rate, so they estimate the realistic case
honestly. This separates the results from the $98$--$99\%$ accuracies reported in
SMOTE-based studies, which train, and sometimes evaluate, on balanced data and
whose numbers therefore describe a different setting. The lower numbers here
reflect a more realistic task rather than weaker models. On the critical side,
undersampling discards information, as $5{,}939$ healthy firms enter each
training fold through only $660$ representatives. With roughly nine firms per
cluster, each representative is estimated from little data and would shift under
a new sample, a trade-off between representativeness and stability that the
repeated cross-validation softens but does not remove. 
The natural benchmark for this trade-off is \textcite{wangliu2021}, whose
framework this thesis leans on. Their best configuration reaches an $F_2$ of
$0.423$ where the best model here reaches $0.550$. This comparison is
encouraging but not exact, as the specifics of the pipelines still differ, for
example in the classifier itself. This is also why the thesis places its
contribution on the feature-importance layer rather than on performance alone.
What the comparison does show is that tree ensembles combined with
cluster-based undersampling, a combination \textcite{wangliu2021} left
unexplored, are worth the attention.